% Source: tex/chapters/ch06/04_ソフトウェアエンジニアリング運用の基礎.tex
\subsubsection{スクリプト化と自動化}

スクリプト化とは、繰り返し行う作業を簡易なプログラムとして記述することである。自動化とは、人が手作業で行っていた処理を、ツールやスクリプトで再現可能に実行することである。運用では、環境構築、設定変更、配備、監視、通知、ログ収集、バックアップ、復旧試験など、多くの活動が自動化の対象になる。

自動化には三つの効果がある。第一に、遅延を減らせる。第二に、手作業のばらつきや設定ミスを減らせる。第三に、障害時の反応を早め、停止時間を短くできる。さらに、自動化された手順は標準化され、運用サービスとして開発者へ提供しやすくなる。

\begin{pointbox}{実務で見るべきこと}
自動化は「便利だから行う」のではない。変更を安全に早く届けるため、同じ環境を何度でも作れるようにするため、障害時に人間の記憶へ依存しないために行うのである。
\end{pointbox}
